% Preamble of the English EGA edition, written for this edition (dedicated to the public domain under CC0 1.0).
% It defines the notation macros, cross-reference commands and numbered environments used by the sources.
\usepackage[a4paper,inner=2.6cm,outer=3.6cm,top=2.8cm,bottom=2.8cm,marginparwidth=2.7cm,marginparsep=0.3cm]{geometry}
\usepackage[T1]{fontenc}
\usepackage[utf8]{inputenc}
\usepackage{lmodern}
\usepackage[english]{babel}
\usepackage{amsmath,amssymb,amsthm,mathrsfs,mathtools}
\usepackage[shortlabels]{enumitem}
\usepackage[all,cmtip]{xy}
\usepackage{tikz-cd}
\usepackage{longtable,multicol}
\usepackage{marginnote}
\usepackage{xstring}
\usepackage[fit]{truncate}
% Running heads: chapter on left pages, section on right pages, in slanted type and cut to the text width.
\makeatletter
\renewcommand{\ps@headings}{%
  \def\@oddfoot{}\def\@evenfoot{}%
  \def\@oddhead{\truncate{\dimexpr\textwidth-3em\relax}{\slshape\rightmark}\hfil\thepage}%
  \def\@evenhead{\thepage\hfil\truncate{\dimexpr\textwidth-3em\relax}{\slshape\leftmark}}%
  \def\chaptermark##1{\markboth{##1}{}}%
  \def\sectionmark##1{\markright{\thesection.\ ##1}}}
\pagestyle{headings}
% The blank left-hand page before a chapter carries no running head.
\renewcommand{\cleardoublepage}{\clearpage\if@twoside\ifodd\c@page\else\hbox{}\thispagestyle{empty}\newpage\fi\fi}
% A bibliography is headed by the section heading of its file, not by a chapter heading of its own.
\renewenvironment{thebibliography}[1]
  {\list{\@biblabel{\@arabic\c@enumiv}}{\settowidth\labelwidth{\@biblabel{#1}}\leftmargin\labelwidth
     \advance\leftmargin\labelsep\@openbib@code\usecounter{enumiv}\let\p@enumiv\@empty
     \renewcommand\theenumiv{\@arabic\c@enumiv}}\sloppy\clubpenalty4000\@clubpenalty\clubpenalty
     \widowpenalty4000\sfcode`\.\@m}
  {\def\@noitemerr{\@latex@warning{Empty `thebibliography' environment}}\endlist}
\makeatother
\usepackage[hidelinks,bookmarksnumbered,hypertexnames=false]{hyperref}% unique link targets for tagged formulas
\usepackage{bookmark}
\setlength{\emergencystretch}{3em}

% ---- notation ---------------------------------------------------------------
\newcommand{\sh}[1]{\mathscr{#1}}
\newcommand{\bb}[1]{\mathbf{#1}}
\newcommand{\bbGamma}{\boldsymbol{\Gamma}}
\newcommand{\cat}[1]{\mathbf{#1}}
\newcommand{\vphi}{\varphi}
\newcommand{\HH}{\mathrm{H}}
\newcommand{\CHH}{\check{\mathrm{H}}}
\newcommand{\RR}{\mathrm{R}}
\newcommand{\LL}{\mathrm{L}}
\newcommand{\red}{\mathrm{red}}
\newcommand{\kres}{\boldsymbol{k}}
\newcommand{\emp}{\varnothing}
\newcommand{\supertilde}{^{\sim}}% written after its argument: M\supertilde is the sheaf associated with M
\newcommand{\isoto}{\xrightarrow{\,\smash{\raisebox{-0.45ex}{\ensuremath{\scriptstyle\sim}}}\,}}
\newcommand{\simto}{\isoto}
\DeclareMathOperator{\Spec}{Spec}
\DeclareMathOperator{\Spf}{Spf}
\DeclareMathOperator{\Proj}{Proj}
\DeclareMathOperator{\Supp}{Supp}
\DeclareMathOperator{\Der}{Der}
\DeclareMathOperator{\codim}{codim}
\DeclareMathOperator{\rg}{rg}
\DeclareMathOperator{\gr}{gr}
\DeclareMathOperator{\Tor}{Tor}
\DeclareMathOperator{\Ext}{Ext}
\DeclareMathOperator{\Ker}{Ker}
\DeclareMathOperator{\Coker}{Coker}
\providecommand{\Hom}{}\renewcommand{\Hom}{\operatorname{Hom}}
\providecommand{\Im}{}\renewcommand{\Im}{\operatorname{Im}}
\newcommand{\shHom}{\mathscr{H}\!\mathit{om}}
\newcommand{\shExt}{\mathscr{E}\!\mathit{xt}}
\newcommand{\shTor}{\mathscr{T}\!\mathit{or}}
\newcommand{\shProj}{\operatorname{\mathbf{Proj}}}
\newcommand{\setmin}{-}% EGA writes the complement of Y in X as X - Y
\newcommand{\C}{C}% a category, as in 0, 8.1.1
\newcommand{\op}{^{0}}% the opposite category C^0
\newcommand{\Set}{\mathrm{Ens}}% the category of sets
\newcommand{\CHom}{\operatorname{Hom}}% the category of functors
\newcommand{\Gr}{\operatorname{Gr}}% associated graded object
\providecommand{\og}{``}
\providecommand{\fg}{''}
\providecommand{\textsection}{\S}

% ---- original pagination ----------------------------------------------------
% \oldpage[volume]{n}: page n of the original printing, shown in the margin. Between paragraphs the note waits for
% the next paragraph to begin, so that it is set beside that paragraph's first line.
\makeatletter
\let\ega@pendingnote\relax
\newcommand{\oldpage}[2][]{%
  \ifvmode
    \gdef\ega@pendingnote{\footnotesize[#1\if\relax\detokenize{#1}\relax\else\,|\,\fi#2]}%
  \else
    \marginnote{\footnotesize[#1\if\relax\detokenize{#1}\relax\else\,|\,\fi#2]}%
  \fi}
\AddToHook{para/begin}{\ifx\ega@pendingnote\relax\else
  \let\ega@note\ega@pendingnote\global\let\ega@pendingnote\relax\marginnote{\ega@note}\fi}
\makeatother

% ---- numbered paragraphs and statements -------------------------------------
% EGA numbers every paragraph and statement explicitly, e.g. (1.2.3); the number is the optional argument,
% and \label after it records that number.
\makeatletter
\newcommand{\ega@setnum}[1]{\def\@currentlabel{#1}\phantomsection}
\newenvironment{env}[1][]{\par\medskip\ega@setnum{#1}\noindent\textbf{(#1)}\enspace\ignorespaces}{\par\medskip}
% \begin{theorem}[19.6.1][Cohen]: number, then an optional name printed after the kind of statement.
\newcommand{\ega@stmt}[3]{%
  \NewDocumentEnvironment{#1}{O{} o}{\par\medskip\ega@setnum{##1}\noindent
    \textbf{(##1)~#2\IfValueT{##2}{ (##2)}.}\enspace#3\ignorespaces}{\par\medskip}}
% \egabib{key}{n} before a bibliography entry printed as plain text "[n] ...": \cite{key} then prints [n] and links here.
\newcommand{\egabib}[2]{\Hy@raisedlink{\hyper@anchorstart{cite.#1}\hyper@anchorend}%
  \protected@write\@auxout{}{\string\bibcite{#1}{#2}}}
% A label that records a given number: used where a label sits at a place whose current number is a neighbour's.
\newcommand{\egalabel}[2]{\begingroup\def\@currentlabel{#2}\label{#1}\endgroup}
\ega@stmt{proposition}{Proposition}{\itshape}
\ega@stmt{theorem}{Theorem}{\itshape}
\ega@stmt{corollary}{Corollary}{\itshape}
\ega@stmt{lemma}{Lemma}{\itshape}
\ega@stmt{definition}{Definition}{}
\ega@stmt{remark}{Remark}{}
\ega@stmt{remarks}{Remarks}{}
\ega@stmt{example}{Example}{}
\ega@stmt{examples}{Examples}{}
\ega@stmt{notation}{Notation}{}
\ega@stmt{scholium}{Scholium}{}
\newenvironment{namedtheorem}[2][]{\par\medskip\ega@setnum{#1}\noindent\textbf{(#1)~#2.}\enspace\itshape\ignorespaces}{\par\medskip}
\newenvironment{namedlemma}[2][]{\par\medskip\ega@setnum{#1}\noindent\textbf{(#1)~#2.}\enspace\itshape\ignorespaces}{\par\medskip}
\makeatother

% The shown text of a link whose target is not in this edition.
\newcommand{\unlinked}[1]{#1}

% A link target the French sources place at numbered statements: \EGAReaderTarget{target name}{label key}.
\newcommand{\EGAReaderTarget}[2]{\ifvmode\else\leavevmode\fi\hypertarget{#1}{}}

% ---- references ----------------------------------------------------------------
% \sref[volume]{key}: the printed number of the labelled paragraph, in parentheses, linked; with a volume prefix
% when the target is in another chapter, e.g. \sref[0]{0.6.7.4} gives (0, 6.7.4).
\newcommand{\sref}[2][]{\hyperref[#2]{(\if\relax\detokenize{#1}\relax\else#1,~\fi\ref*{#2})}}
% \eref{key}: a numbered formula, e.g. (11.2.1.1).
\newcommand{\eref}[1]{\hyperref[#1]{(\ref*{#1})}}
% \erratum[volume]{reference}: a reference as corrected in the errata of that volume (Err_IV).
\newcommand{\erratum}[2][]{#2~[Err$_{\mathrm{#1}}$]}
